\documentclass[10pt,aspectratio=169]{beamer}
\usepackage[portuguese]{babel}
\usepackage{fontspec,unicode-math}
\setmainfont{Noto Sans}
\setsansfont{Noto Sans}
\setmathfont{STIX Math}
\usepackage{amsmath}
\usepackage{xcolor}
\definecolor{InsperRed}{HTML}{a01924}
\definecolor{Ink}{HTML}{232529}
\setbeamercolor{frametitle}{fg=Ink}
\setbeamercolor{structure}{fg=InsperRed}
\setbeamercolor{normal text}{fg=Ink,bg=white}
\setbeamerfont{frametitle}{size=\large,series=\bfseries}
\setbeamerfont{title}{size=\LARGE,series=\bfseries}
\setbeamertemplate{navigation symbols}{}
\setbeamertemplate{itemize item}{\color{InsperRed}$\bullet$}
\setbeamertemplate{footline}{%
  \hbox{\begin{beamercolorbox}[wd=.84\paperwidth,ht=2.7ex,dp=1.2ex,leftskip=.7cm]{normal text}\scriptsize Inteligência Artificial para Decisão · Insper · 2027-1\end{beamercolorbox}%
  \begin{beamercolorbox}[wd=.16\paperwidth,ht=2.7ex,dp=1.2ex,rightskip=.7cm plus1fil]{normal text}\hfill\scriptsize\insertframenumber\, /\, \inserttotalframenumber\end{beamercolorbox}}}
\setbeamersize{text margin left=12mm,text margin right=12mm}
\begin{document}
\begin{frame}[t]{}
\vspace{14mm}\textcolor{InsperRed}{\bfseries AULA 11}\par\vspace{4mm}{\LARGE\bfseries Workflows e agentes de IA}\par\vspace{6mm}{\large Operações empresariais}\par\vspace{12mm}{\color{gray}Paulo C. Marques F. · 2027-1}
\end{frame}

\begin{frame}[t]{Objetivos e verificação}
\small
\begin{itemize}
\item Distinguir workflow predefinido de comportamento agentivo, modelar estados e transições e instrumentar chamadas de ferramentas em ambiente simulado.
\item Resolver manualmente os exemplos do capítulo, identificando hipóteses e consequências.
\item Especificar a aplicação no Codex gráfico e confrontar ao menos um resultado com evidências.
\end{itemize}
\end{frame}

\begin{frame}[t]{Organização do encontro}
\small
\begin{itemize}
\item Apresentação dos fundamentos e resolução de problemas quantitativos no papel.
\item Laboratório supervisionado e testes produzidos por prompts na interface gráfica do Codex, sem terminal.
\item Atividade individual com dados diferentes e depósito no GitHub antes do quiz seguinte.
\end{itemize}
\end{frame}

\begin{frame}[t]{Fluxos determinísticos e decisões delegadas}
\small
\begin{itemize}
\item Um workflow especifica antecipadamente etapas, estados, condições e rotas possíveis.
\item Um agente pode decidir dinamicamente qual ferramenta utilizar e quais passos seguir, dentro de permissões definidas.
\item A distinção não é binária: uma aplicação pode ter partes determinísticas e trechos delegados ao modelo.
\item Em todos os casos, efeitos sobre sistemas externos devem permanecer sujeitos a regras de autorização.
\end{itemize}
\end{frame}

\begin{frame}[t]{Fluxos determinísticos e decisões delegadas — continuação}
\small
\begin{itemize}
\item O estado pode incluir identificador do pedido, data de entrega, motivo da solicitação, status da aprovação e resultado de cada ferramenta.
\item Uma transição é uma função , em que é ação proposta e é retorno observado.
\item Os estados e transições precisam ser auditáveis.
\end{itemize}
\end{frame}

\begin{frame}[t]{Exemplo de política de devolução}
\small
\begin{itemize}
\item Suponha uma política didática: devoluções por arrependimento podem ser solicitadas até 7 dias após recebimento; depois desse prazo, o sistema encaminha para análise humana.
\item Para um pedido entregue há 5 dias, a condição temporal é satisfeita; entregue há 9 dias, não é.
\item O prazo não deve ser inventado pelo LLM, mas consultado na política efetivamente fornecida.
\item Situações de defeito podem exigir outra rota; o exemplo não pretende representar a legislação aplicável.
\end{itemize}
\end{frame}

\begin{frame}[t]{Exemplo de política de devolução — continuação}
\small
\begin{itemize}
\item A ferramenta consultar\_pedido retorna data e situação.
\item A ferramenta consultar\_politica retorna a regra.
\item A aplicação pode propor abrir\_solicitacao, mas não deve executar efetuar\_estorno sem aprovação explícita e verificável.
\end{itemize}
\end{frame}

\begin{frame}[t]{Controle de ações e idempotência}
\small
\begin{itemize}
\item Um teste importante verifica se a repetição de uma solicitação não executa duas vezes o mesmo efeito.
\item Esse princípio é conhecido como idempotência operacional.
\item Um identificador único de requisição e um registro de status permitem rejeitar operações duplicadas.
\item Limitar o número de passos do agente, registrar timeouts e exigir condição de parada previne loops.
\end{itemize}
\end{frame}

\begin{frame}[t]{Rastros e testes}
\small
\begin{itemize}
\item O laboratório produzirá rastros com estado inicial, ferramentas chamadas, argumentos, retornos, condição avaliada e estado final.
\item Um bom teste inclui pedidos elegíveis, inelegíveis, dados ausentes, políticas conflitantes e solicitações repetidas.
\item Um agente com resposta textual correta mas com uma operação não autorizada executada falhou.
\item A avaliação deve distinguir trajetória, resultado e respeito às permissões.
\end{itemize}
\end{frame}

\begin{frame}[t]{Exercício para resolver no papel 1}
\small
\small\textbf{11.1 O sétimo dia}\par\vspace{4mm}Uma regra didática permite devolução até 7 dias corridos após a entrega, incluindo o próprio sétimo dia. Classifique pedidos recebidos há 6, 7 e 8 dias; formule a guarda booleana que decide elegibilidade.
\end{frame}

\begin{frame}[t]{Exercício para resolver no papel 2}
\small
\small\textbf{11.2 Um grafo com caminho proibido}\par\vspace{4mm}Os estados são recebido, identificado, verificado, aprovado e concluído. O grafo admite aresta direta recebido concluído. Qual controle falta? Descreva uma condição que impossibilite concluir sem identificar e verificar o pedido.
\end{frame}

\begin{frame}[t]{Laboratório e caso aplicado}
\small
\begin{itemize}
\item Área de aplicação: Operações empresariais.
\item Desenhar à mão grafo de estados recebido, identificado, verificado, aguardando-aprovacao, concluido e escalado, com duas guardas.
\item Pedir ao Codex simulador visual com ferramentas de leitura, transições determinísticas e log de estado anterior/ação/resposta/estado posterior.
\item Executar cinco cenários: elegível, fora do prazo, pedido inexistente, indisponibilidade e chamada duplicada.
\end{itemize}
\end{frame}

\begin{frame}[t]{Prompt de desenvolvimento — exemplo}
\small
\small\textbf{Instrução fornecida ao Codex:}\par\vspace{5mm}Crie simulador local de trocas/devoluções com estados, guardas e log auditável. Todos os efeitos devem ser fictícios. Implemente interfaces gráficas de inspeção, testes de erro, interrupção e retomada; ações de escrita exigem validação por código independente de modelos.
\end{frame}

\begin{frame}[t]{Critérios de aceitação e validação}
\small
\begin{itemize}
\item Estado final deriva de eventos persistidos; nenhuma guarda ignorada; ações duplicadas tratadas idempotentemente; falhas geram estado de exceção; rastros podem ser reconstituídos por pedido.
\item Inspecionar argumentos, saídas e rastros das ferramentas quando houver chamadas externas.
\item Comparar uma saída com a resolução manual do exemplo conceitual.
\end{itemize}
\end{frame}

\begin{frame}[t]{Atividade individual em casa}
\small
\begin{itemize}
\item Simulador alternativo de mudança de endereço pós-compra, com oito pedidos fictícios, política própria e dois caminhos de exceção.
\item Resolver as questões de interpretação, com cálculos e evidência documental.
\item Depositar no GitHub os arquivos solicitados; a entrega habilita o quiz, sem valer nota.
\end{itemize}
\end{frame}

\begin{frame}[t]{Preparação para a avaliação escrita}
\small
\begin{itemize}
\item Exercícios do caderno: 11.1 O sétimo dia; 11.2 Um grafo com caminho proibido; 11.3 Escolher ou seguir uma regra?.
\item Formalizar o raciocínio e interpretar resultados; não será exigida escrita manual de código.
\item No quiz seguinte serão avaliados conceitos desta aula e da atividade feita em casa.
\end{itemize}
\end{frame}

\end{document}
